\epanalipsi{Α6}

\tmima{Λεξικό του κεφαλαίου}

\renewcommand{\arraystretch}{1.5}
\begin{tabular*}{\linewidth}{@{\extracolsep{\fill}} >{\raggedright\arraybackslash}p{0.28\linewidth} >{\raggedright\arraybackslash}p{0.62\linewidth}}
\toprule
\textbf{Λέξη ή σύμβολο} & \textbf{Με μία πρόταση}\\
\midrule
ισότητα $\alpha=\beta$ & Τα δύο μέλη έχουν την ίδια τιμή, όπως μια ζυγαριά που ισορροπεί.\\
ιδιότητες της ισότητας & Αν $\alpha=\beta$, τότε $\alpha\pm\gamma=\beta\pm\gamma$, $\alpha\gamma=\beta\gamma$ και $\alpha:\gamma=\beta:\gamma$ (με $\gamma\neq0$).\\
εξίσωση & Ισότητα με έναν άγνωστο αριθμό, π.χ. $2x+6=14$.\\
λύση ή ρίζα & Ο αριθμός που, στη θέση του αγνώστου, κάνει την ισότητα αληθινή.\\
επίλυση, επαλήθευση & Η εύρεση της λύσης· ο έλεγχος της λύσης στην αρχική εξίσωση.\\
ισοδύναμες εξισώσεις & Εξισώσεις με τις ίδιες λύσεις.\\
αδύνατη, ταυτότητα & Καμία λύση (π.χ. $0\cdot x=2$)· όλοι οι αριθμοί είναι λύση ($0\cdot x=0$).\\
$x=\dfrac{\gamma-\beta}{\alpha}$ & Η λύση της $\alpha x+\beta=\gamma$, με $\alpha\neq0$.\\
\bottomrule
\end{tabular*}

\tmima{Ασκήσεις επανάληψης}

\ex Αν $\alpha=\beta$, συμπλήρωσε με το σωστό: $\alpha+7=\dots$, \quad $\alpha-2=\dots$, \quad $5\alpha=\dots$, \quad $\alpha:4=\dots$

\ex Εξέτασε αν το $4$ είναι λύση των εξισώσεων: α) $x+9=13$ \quad β) $3x=15$ \quad γ) $20-x=16$ \quad δ) $2x+1=10$.

\ex Λύσε: α) $x+19=40$ \quad β) $x-2{,}5=4$ \quad γ) $30-x=12$ \quad δ) $6x=42$ \quad ε) $x:7=3$ \quad στ) $56:x=8$.

\ex Λύσε και επαλήθευσε: α) $-4x=28$ \quad β) $x+(-3)=-10$ \quad γ) $5-x=-2$.

\ex Λύσε: α) $5x+3=28$ \quad β) $4x-9=15$ \quad γ) $2-3x=14$ \quad δ) $0{,}5x+2=3{,}5$.

\ex Λύσε: α) $\dfrac{2}{3}x+1=5$ \quad β) $\dfrac{x}{4}-2=1$.

\ex Λύσε: $2(x+5)-3(x-1)=10$.

\ex Λύσε: $3(2x-1)-(x+4)=12$.

\ex Εξέτασε αν η εξίσωση $4(x+2)=4x+8$ είναι αδύνατη ή ταυτότητα και αν η $3x+1=3x-4$ είναι αδύνατη ή ταυτότητα.

\ex Γράψε εξίσωση και λύσε: το τριπλάσιο ενός αριθμού αυξημένο κατά $8$ κάνει $29$.

\ex Η Αγγελική αγόρασε $3$ ίδια τετράδια και μια κασετίνα των $6$ €. Πλήρωσε συνολικά $15$ €. Πόσο κοστίζει κάθε τετράδιο;

\ex Η Ζωή έχει $x$ ευρώ και ο Νίκος $7$ ευρώ περισσότερα. Μαζί έχουν $33$ ευρώ. Πόσα έχει η Ζωή;

\ex Ένα ορθογώνιο έχει πλάτος $8$ cm και περίμετρο $54$ cm. Βρες το μήκος του με εξίσωση.

\ex Ο παππούς Γιάννης λέει: «Αν διπλασιάσεις την ηλικία μου και αφαιρέσεις $20$, βρίσκεις $130$». Πόσο ετών είναι;

\ex \textbf{Σωστό ή λάθος (ισότητα).} α) Αν $\alpha=\beta$, τότε $\alpha-5=\beta-5$. \quad β) Αν $3\alpha=3\beta$, τότε $\alpha=\beta$. \quad γ) Η $0\cdot x=5$ είναι ταυτότητα. \quad δ) Αν $\alpha\cdot0=\beta\cdot0$, τότε $\alpha=\beta$.

\ex \textbf{Σωστό ή λάθος (εξισώσεις).} α) Η $2x=2$ έχει λύση το $2$. \quad β) Η $x+1=x+2$ είναι αδύνατη. \quad γ) Η λύση της $\alpha x+\beta=\gamma$ ($\alpha\neq0$) είναι $x=\dfrac{\gamma-\beta}{\alpha}$. \quad δ) Οι $x+3=7$ και $2x=8$ είναι ισοδύναμες.

\tmima{Ο χάρτης της επανάληψης}

Μετά τις λύσεις, δες σε ποιες ασκήσεις δυσκολεύτηκες. Ο πίνακας σου λέει ποια παράγραφο να ξαναδιαβάσεις.

\renewcommand{\arraystretch}{1.5}
\begin{tabular*}{\linewidth}{@{\extracolsep{\fill}} >{\raggedright\arraybackslash}p{0.38\linewidth} >{\raggedright\arraybackslash}p{0.52\linewidth}}
\toprule
\textbf{Αν δυσκολεύτηκες στις} & \textbf{ξαναδιάβασε τις παραγράφους}\\
\midrule
ΕΑ6.1, ΕΑ6.14 & 6.1 (ιδιότητες της ισότητας)\\
ΕΑ6.2, ΕΑ6.3, ΕΑ6.4, ΕΑ6.9 & 6.1 (εξίσωση, βασικές μορφές, ειδικές περιπτώσεις)\\
ΕΑ6.5, ΕΑ6.6, ΕΑ6.7, ΕΑ6.8 & 6.2 (η μέθοδος $\alpha x+\beta=\gamma$)\\
ΕΑ6.10, ΕΑ6.11, ΕΑ6.12, ΕΑ6.13, ΕΑ6.15 & 6.2 (προβλήματα με εξίσωση)\\
ΕΑ6.16 & 6.1 και 6.2\\
\bottomrule
\end{tabular*}
